\section{Recuperatorio 2023 Q1}

% =====================================================================
\subsection{Ejercicio 1 — Tiempo de ataque (fórmula de Anderson)}
% =====================================================================

\begin{consigna}
¿Cuál es el tiempo que un atacante necesita probar passwords en un sistema de autenticación
para que la probabilidad de adivinar sea $1/100$ si con un TPU se pueden probar $700$ claves por
segundo y la clave tiene $12$ bits de longitud?
\end{consigna}

\paragraph{Temas.} Clase 7 — Autenticación (\S Fórmula de Anderson: subir la complejidad).

\paragraph{Qué necesitás saber.}
La \textbf{fórmula de Anderson} estima la probabilidad de que un ataque de fuerza bruta acierte
una clave:
\begin{equation}
  P \;=\; \frac{T\cdot G}{N},
\end{equation}
donde $P$ = probabilidad de adivinar, $T$ = tiempo dedicado a las pruebas (en segundos),
$G$ = pruebas por segundo y $N$ = tamaño del espacio de claves. Despejando el tiempo:
$T \le P\,N/G$. Cuidado con dos cosas: (i) la ``longitud en bits'' fija el espacio como
$N = 2^{\,\text{bits}}$ (no hay que elevar un alfabeto); (ii) la unidad de $T$ es segundos.

\paragraph{Notas y conexiones.}
Es el mismo despeje de los tres ejercicios resueltos del informe (Clase 7): el caso ``calcular
el tiempo'' usa $T \le P\,N/G$ y el caso ``longitud mínima'' usa $N \ge T\,G/P$ seguido de
$L \ge \log_{\text{base}} N$. Acá la clave está dada \emph{en bits}, así que el espacio es
directamente una potencia de $2$ y no hace falta pasar por logaritmos.

\paragraph{Resolución.}
El espacio de claves de $12$ bits es
\begin{equation}
  N = 2^{12} = 4096 .
\end{equation}
Con $G = 700$ claves/s y $P = 1/100 = 0{,}01$:
\begin{align*}
  T &\le \frac{P\,N}{G} = \frac{0{,}01 \cdot 4096}{700}
       = \frac{40{,}96}{700} \approx 0{,}0585\ \text{s}.
\end{align*}
El atacante necesita probar durante \textbf{$\approx 0{,}0585$ segundos} (unos $59$ ms) para
alcanzar una probabilidad de éxito de $1/100$. El número es ridículamente chico: un espacio de
$2^{12}$ claves es \textbf{trivialmente rompible} (de hecho, probando las $4096$ claves a
$700$/s se agota el espacio entero en $\approx 5{,}85$ s, garantizando el acierto). La conclusión
de fondo es que $12$ bits no alcanzan ni de lejos como longitud de clave.

% =====================================================================
\subsection{Ejercicio 2 — Secreto compartido de Shamir $(3,4)$ en $\mathbb{Z}_7$}
% =====================================================================

\begin{consigna}
Considerar un esquema de secreto compartido de Shamir $(3,4)$ en $\mathbb{Z}_7$, con el conjunto
de $4$ sombras $C = \{(1,4),(2,1),(3,1),(5,3)\}$.
\begin{enumerate}
  \item[a)] Recuperar el secreto.
  \item[b)] Indicar el polinomio utilizado para obtener el conjunto de sombras.
\end{enumerate}
\end{consigna}

\paragraph{Temas.} Clase 8 — Principios de Diseño y Vulnerabilidades (\S El secreto compartido de
Shamir, en términos de entropía); base aritmética de la Clase 9 — Flujo de Información (entropía).

\paragraph{Qué necesitás saber.}
En un esquema de Shamir $(T,N)$ hay $N$ sombras y se necesitan \textbf{al menos $T$} para
recuperar el secreto. El reparto se hace con un polinomio de grado $T-1$ sobre un cuerpo finito
$\mathbb{Z}_p$:
\begin{equation}
  f(x) = a_0 + a_1 x + a_2 x^2 + \dots + a_{T-1}x^{\,T-1} \pmod p,
\end{equation}
donde el \textbf{secreto es $S = a_0 = f(0)$} y cada sombra es un par $(x_i, f(x_i))$. Con $T$
pares cualesquiera se reconstruye unívocamente el polinomio (interpolación de Lagrange) y se
evalúa en $0$. Acá $T=3$, $N=4$, $p=7$, así que el polinomio es de grado $2$ y alcanza con
\textbf{$3$ de las $4$ sombras}. Toda la aritmética es módulo $7$ (los inversos se calculan
con el inverso modular).

\paragraph{Notas y conexiones.}
El informe (Clase 8) justifica la validez del esquema \emph{con entropía}: con menos de $T$
sombras la entropía del secreto \textbf{no cambia} ($H(S\mid S_1)=H(S\mid S_1,S_2)=H(S)$, no hay
flujo de información), y recién con $T$ sombras $H(S\mid S_1,S_2,S_3)=0$ (se conoce el secreto).
Es el mismo concepto de flujo de información de la Clase 9: conocer ``demasiado poco'' no reduce
la incertidumbre. La sombra sobrante $(5,3)$ sirve para \emph{verificar} (todo subconjunto de $3$
debe dar el mismo secreto) y para tolerar la pérdida de una sombra.

\paragraph{Resolución.}
\textbf{(a) Recuperar el secreto.} Usamos interpolación de Lagrange evaluada en $x=0$ con tres
sombras, por ejemplo $(1,4),(2,1),(3,1)$:
\begin{equation}
  S = f(0) = \sum_{i} y_i \prod_{j\neq i} \frac{0 - x_j}{x_i - x_j} \pmod 7 .
\end{equation}
Calculando los coeficientes de Lagrange en $0$ (todo mod $7$, con inversos
$2^{-1}\equiv 4,\ 6^{-1}\equiv 6,\ 5^{-1}\equiv 3$):
\begin{align*}
  L_1(0) &= \frac{(0-2)(0-3)}{(1-2)(1-3)} = \frac{6}{2} \equiv 6\cdot 4 \equiv 3,\\
  L_2(0) &= \frac{(0-1)(0-3)}{(2-1)(2-3)} = \frac{3}{-1} \equiv 3\cdot 6 \equiv 4,\\
  L_3(0) &= \frac{(0-1)(0-2)}{(3-1)(3-2)} = \frac{2}{2} \equiv 1.
\end{align*}
\begin{align*}
  S &= 4\cdot L_1(0) + 1\cdot L_2(0) + 1\cdot L_3(0)
     = 4\cdot 3 + 1\cdot 4 + 1\cdot 1 \\
    &= 12 + 4 + 1 = 17 \equiv \boxed{3} \pmod 7 .
\end{align*}
El secreto es $\mathbf{S = 3}$. (Se verifica con cualquier otro subconjunto de $3$ sombras: todos
dan $3$, lo que confirma que las cuatro sombras son consistentes.)

\medskip
\textbf{(b) Polinomio.} Buscamos $f(x) = a_0 + a_1 x + a_2 x^2 \pmod 7$ con $a_0 = S = 3$.
Imponiendo dos sombras más, $f(1)=4$ y $f(2)=1$:
\begin{align*}
  f(1) &= 3 + a_1 + a_2 \equiv 4 &&\Rightarrow\; a_1 + a_2 \equiv 1 \pmod 7,\\
  f(2) &= 3 + 2a_1 + 4a_2 \equiv 1 &&\Rightarrow\; 2a_1 + 4a_2 \equiv -2 \equiv 5 \pmod 7.
\end{align*}
De la primera, $a_1 \equiv 1 - a_2$. Sustituyendo:
$2(1-a_2) + 4a_2 \equiv 5 \Rightarrow 2 + 2a_2 \equiv 5 \Rightarrow 2a_2 \equiv 3
\Rightarrow a_2 \equiv 3\cdot 4 \equiv 12 \equiv 5 \pmod 7$, y entonces
$a_1 \equiv 1 - 5 \equiv -4 \equiv 3 \pmod 7$. El polinomio es
\begin{equation}
  \boxed{\,f(x) = 3 + 3x + 5x^2 \pmod 7\,}.
\end{equation}
Verificación de las cuatro sombras: $f(1)=11\equiv4$, $f(2)=25\equiv4{+}\dots$;
explícitamente $f(1)\equiv 4$, $f(2)\equiv 1$, $f(3) = 3+9+45 = 57 \equiv 1$,
$f(5) = 3+15+125 = 143 \equiv 3 \pmod 7$. Las cuatro coinciden con $C$.

% =====================================================================
\subsection{Ejercicio 3 — Verdadero o falso (justificando)}
% =====================================================================

\begin{consigna}
Responder verdadero o falso, justificando o corrigiendo como fuere apropiado:
\begin{enumerate}
  \item[a)] En sistemas de control de acceso que incluyen agrupamientos de usuarios, la
  resolución de conflictos en la asignación de los permisos de GRANT-ALL implica darle a todos
  los usuarios posibles los permisos.
  \item[b)] Los esquemas de autenticación por varios factores ofrecen esquemas más seguros de
  romper debido a que se requiere comprometer de manera simultanea varios canales de
  autenticación.
  \item[c)] En Bell-Lapadula la condición de seguridad simple (CSS) implica que un usuario puede
  leer un objeto si y solo si su nivel es igual o mayor que el del objeto.
  \item[d)] En un sistema sensible de generación de claves simétricas, si la entropía de esa clave
  es $256$ bits, pero la entropía de esa clave se reduce a $200$ bits si se registra el sonido del
  disipador de hardware, ¿Hay flujo de información de una variable a la otra?
\end{enumerate}
\end{consigna}

\paragraph{Temas.}
(a) Clase 6 — Políticas y Control de Acceso (\S Conflictos en ACL);
(b) Clase 7 — Autenticación (\S MFA);
(c) Clase 6 — Políticas y Control de Acceso (\S Bell--LaPadula, seguridad simple);
(d) Clase 9 — Flujo de Información y Malware (\S Entropía y flujo de información / canal lateral).

\paragraph{Qué necesitás saber.}
\begin{itemize}
  \item \textbf{(a) Resolución de conflictos en ACL.} Cuando varias entradas (de un sujeto y de
  los grupos a los que pertenece) le dan permisos distintos sobre un objeto, hay varias políticas:
  \emph{Permitir} (unión: el acceso si \emph{alguna} entrada lo da), \emph{Grant-All}
  (\textbf{intersección}: denegar si \emph{alguna} entrada lo deniega $\Rightarrow$ exige que
  \emph{todas} den el permiso) y \emph{First-Match}. Grant-All es la regla \textbf{más restrictiva},
  no la más permisiva.
  \item \textbf{(b) MFA.} Multi-factor combina factores de \emph{distinta naturaleza} (algo que
  conozco / tengo / soy); para romperlo hay que comprometer varios \emph{assets} distintos a la vez.
  \item \textbf{(c) Bell--LaPadula, seguridad simple (\emph{read-down}).} El sujeto $s$ puede
  \emph{leer} el objeto $o$ si y solo si $L(s) \dom L(o)$, es decir su nivel \textbf{domina}
  (es igual o superior) al del objeto: ``read down''.
  \item \textbf{(d) Flujo de información (entropía).} Hay flujo de $x$ a $y$ si la incertidumbre
  sobre $x$ \emph{disminuye} al observar $y$ ($H(x\mid y) < H(x)$). \textbf{Reducir entropía es
  flujo de información.}
\end{itemize}

\paragraph{Notas y conexiones.}
El ítem (d) es el ejercicio canónico del disipador (señalado por el índice como el clásico de este
recuperatorio): la entropía de la clave baja de $256$ a $200$ bits al registrar el sonido del
disipador, o sea se filtran $56$ bits. Es una doble lectura: el sonido del disipador es una
\textbf{vía no diseñada para transmitir} $\Rightarrow$ \emph{canal encubierto}; explotar esa
emisión física del hardware para sacar la clave es un \textbf{ataque de canal lateral}
(``dos caras de la misma moneda''). El (c) es la trampa de direcciones de BLP: confundir
\emph{read-down} (lectura: dominás al objeto) con \emph{write-up} es el error clásico; aquí el
enunciado describe correctamente la lectura. El (a) es el típico ejercicio de conflictos de ACL
del informe (Clase 6), donde Grant-All resulta en la \emph{intersección} $\{w\}$, no en todos los
permisos.

\paragraph{Resolución.}

\textbf{(a) FALSO.} Grant-All es exactamente la política \emph{opuesta}: resuelve el conflicto por
\textbf{intersección} de los permisos, es decir \emph{deniega} salvo que \textbf{todas} las
entradas que aplican al usuario (la suya y las de sus grupos) concedan el permiso. Es la regla
\emph{más restrictiva}, no la que da ``todos los permisos posibles''. (Dar el permiso si
\emph{alguna} entrada lo concede es la política de \emph{unión/permitir}, que es la otra.) En el
ejemplo resuelto del informe, con $\mathrm{ACL}(o)=\{(G,rw),(s_1,w),(\{s_1,s_2\},wx)\}$ y
$s_1\in G$, Grant-All da $rw \cap w \cap wx = \{w\}$ (solo lo que está en \emph{todas}), mientras
que la unión daría $\{r,w,x\}$.

\medskip
\textbf{(b) VERDADERO} (con un matiz). MFA es más robusto porque exige comprometer \textbf{varios
factores simultáneamente}, y la idea de diseño es que sean de \textbf{naturaleza distinta} (algo
que conozco, algo que tengo, algo que soy), de modo que romperlos implica vulnerar
\emph{assets/canales} diferentes. El matiz que conviene aclarar: el beneficio depende de que los
factores sean \emph{independientes}; si se concentran en un mismo dispositivo (el celular reúne
``tengo'' + ``soy''), clonarlo derriba ambos a la vez y la ganancia se diluye. La afirmación, tal
como está, es correcta.

\medskip
\textbf{(c) VERDADERO.} La condición de seguridad simple (\emph{simple security}, ``read down'')
de Bell--LaPadula dice que $s$ puede leer $o$ si y solo si $L(s) \dom L(o)$, o sea cuando el nivel
del sujeto es \textbf{igual o mayor (domina)} al del objeto. El enunciado lo describe
correctamente.
\begin{equation*}
  s \text{ puede leer } o \iff L(s) \dom L(o) \qquad (\text{read down}).
\end{equation*}
Conviene precisar: con compartimentos, ``igual o mayor'' es la relación de \emph{dominancia}
$(L,C)\dom(L',C') \iff L'\le L \wedge C'\subseteq C$ (no un simple orden total), y debe sumarse el
permiso discrecional de lectura. La complementaria es la propiedad-$\star$ (\emph{write up}),
$L(o)\dom L(s)$, que el enunciado no menciona.

\medskip
\textbf{(d) SÍ, hay flujo de información.} El criterio formal es que hay flujo de $x$ (la clave) a
$y$ (la observación del sonido del disipador) cuando la incertidumbre sobre $x$ disminuye al
observar $y$:
\begin{equation*}
  H(x \mid y) < H(x).
\end{equation*}
Aquí la entropía de la clave pasa de $H(x) = 256$ bits a $H(x\mid y) = 200$ bits, es decir
\textbf{baja en $56$ bits}, que es la información que se filtra de la clave hacia el atacante:
\begin{equation*}
  H(x) - H(x\mid y) = 256 - 200 = 56 \text{ bits filtrados} \;>\; 0.
\end{equation*}
Como la incertidumbre disminuye, \textbf{hay flujo de información} de la clave al canal acústico.
Conceptualmente: el sonido del disipador es una \emph{vía no diseñada para transmitir} (canal
encubierto), y aprovechar esa emisión física del hardware para inferir bits de la clave es un
\emph{ataque de canal lateral}. El error que penaliza sería no reconocer que reducir entropía
\emph{es} flujo de información.
